 \chapter{A Visão MedX (Medical Experience): A Perspectiva Técnica}

O Medical Experience (MedX) é o parâmetro de design que governa cada linha de código do IntraClinica Doc Expander. Diferentemente da UX (User Experience) genérica, o MedX considera as condições específicas e complexas do ambiente clínico, abrangendo requisitos tais como iluminação adequada, necessidade de assepsia estrita, e valor do contato visual.

\section{O Diagnóstico da Situação Clínica}

A concepção do IntraClinica Doc Expander se baseou em duas condições clínicas fundamentais observadas na prática tradicional:

\subsection{Hemorragia Financeira no Almoxarifado}
Foi identificado que o sumistro de insumos críticos está frequentemente ausente, resultando em um controle manual ineficiente e consequentemente em perdas financeiras invisíveis mas constantes. O IntraClinica Doc Expander funciona como um estancador dessa sangria financeira, monitorando e controlando de forma automatizada o fluxo dos insumos.

\subsection{Complexidade Paralisante}
Os softwares disponíveis no mercado oferecem uma variedade vastísima de funcionalidades que geralmente não são utilizadas, porém a sobrecarga dessas funções acarreta em uma interface complicada e demorada. Essa obesidade de recursos pode causar fadiga ao profissional de saúde, resultando em um preenchimento abreviado do prontuário apenas com o mínimo necessário para o faturamento. O IntraClinica Doc Expander busca reduzir a complexidade na interface, priorizando a função sobre a quantidade de recursos disponíveis.

\section{Interface Minimalista e a Regra da Uma Mão}

O IntraClinica Doc Expander se encontra desenvolvido de acordo com a filosofia de que ``se o sistema precisa de treinamento, ele falhou''. Para o médico, isso significa uma interface otimizada para tablets, possibilitando navegação plena e registro de anamnese utilizando apenas uma mão. Isso garante que a tecnologia permaneça em segundo plano, preservando a empatia na relação médico-paciente. Além disso, o IntraClinica Doc Expander foi projetado para minimizar o impacto nos hábitos de trabalho existentes dos profissionais de saúde, garantindo sua eficiência ao longo do tempo.

O IntraClinica Doc Expander foi desenvolvido com base em práticas clínicas e técnicas inovadoras, visando aumentar a eficiência dos profissionais de saúde enquanto mantém uma interface intuitiva e confortável. Além disso, o IntraClinica Doc Expander foi projetado para ser adaptável às necessidades específicas do setor de saúde, oferecendo recursos avançados tais como análise de imagens em alta definição e integração com sistemas eletrônicos de saúde.